\documentclass[10pt,a4paper]{amsart}
\usepackage[margin=19mm]{geometry}
\usepackage{amsmath,amssymb,mathtools}
\usepackage[scr=rsfs,cal=euler]{mathalfa}
\usepackage{xfrac}
\usepackage[unicode,hidelinks,bookmarks=false]{hyperref}
\newcommand{\ie}{i.e.\ }
\newcommand{\cf}{cf.\ }
\newcommand{\resp}{respectively\ }
\newcommand{\Subsection}{\subsection}
\providecommand{\nobreakdash}{\nobreak\hbox{-}}
\setlength{\emergencystretch}{2em}
\multlinegap0pt
\usepackage{amssymb}
\usepackage{mathtools}
\usepackage[shortlabels]{enumitem}
\usepackage{xcolor}
\newtheorem{proposition}{Proposition}
\newcommand{\Span}{\operatorname{span}}
\title{Counterexamples to Conjectures of Wehlau on Noether Numbers}
\author{Muhammad Fazeel Anwar}
\date{Source: arXiv:2607.18585v2 (CC BY 4.0)}
\begin{document}
\maketitle

\noindent\emph{Source and licence.} Counterexamples to Conjectures of Wehlau on Noether Numbers. Version arXiv:2607.18585v2 (CC BY 4.0), distributed by arXiv under the Creative Commons Attribution 4.0 International licence, http://creativecommons.org/licenses/by/4.0/, as recorded in the arXiv licence metadata for this exact version and shown on the abstract page https://arxiv.org/abs/2607.18585v2. Full licence notice: \texttt{licenses/arxiv-2607.18585-CC-BY-4.0.txt}.\par\smallskip

\noindent\textit{Editorial scope.} Counted unit: the article's proposition prop2 (source0.tex lines 249--253). The article's complete original proof of it is delivered with it (source0.tex lines 255--376, one \texttt{proof} environment of 122 lines). No mathematics is written, repaired or re-derived: the statement and the proof are the article's own lines, in the article's own notation, and no retained line is substituted or re-flowed.  No further source-local context is needed: every cross-reference of every retained line resolves inside the two delivered spans.  Nothing was omitted from any retained span, so the package carries no omission record.  No retained line contains a citation, so the wrapper carries no bibliography and no bibliography entry.  No retained line contains a figure environment, an image-inclusion command or drawing code, so no image is delivered and the package declares no asset dependency.  Editorial formatting only, in addition: an AMS \texttt{amsart} preamble that restates the article's own declarations and macros used by the retained lines, namely the environments \texttt{proposition} on separate counters, together with the standard packages those lines need.  The retained environments print in this document as Proposition 1 in order of appearance, whereas the article prints the same items with the numbering of its own full document.  The complete native source of the article as distributed in the arXiv e-print for this exact version is bundled unchanged as \texttt{sources/205660/source0.tex}.
	\begin{proposition}\label{prop2}
		Let $S=\mathbb F _2[x_0,x_1,x_2,x_3,x_4]^G$. Then \(S\) is generated by the homogeneous invariants \(\{f_1,\ldots,f_9\}\), with $\deg(f_1),\ldots,\deg(f_9)	= 1,1,2,3,4,4,4,4,6$. The invariants are $ f_1=x_0, \qquad f_2=x_3, \qquad f_3=x_0x_1+x_1^2, \qquad f_4= x_0^2x_4+x_0x_1x_3+x_0x_4^2+x_1x_3^2, \qquad f_5= x_0^3x_2+x_0^2x_1^2+x_0^2x_1x_4+x_0^2x_2^2+x_0x_1^3+x_0x_1^2x_3+x_0x_1^2x_4+x_1^3x_3, \qquad f_6= x_0^3x_4+x_0x_1^2x_3+x_0x_3^2x_4+x_1x_3^3+x_3^2x_4^2+x_4^4, \qquad f_7= x_0^2x_1x_2+x_0^2x_2^2+x_0x_1^2x_2+x_0x_1x_2^2+x_1^2x_2^2+x_2^4, \qquad f_8=x_0^2x_1x_4+x_0^2x_2x_3+x_0x_1^2x_4+x_0x_1x_3x_4+x_0x_1x_4^2+x_0x_2^2x_3+x_1^2x_3x_4+x_1^2x_4^2, \qquad f_9=x_0^4x_1x_4+x_0^4x_2x_4+x_0^3x_1x_2x_3+x_0^3x_1x_3x_4+x_0^3x_2^2x_4+x_0^3x_2x_4^2+x_0^3x_3^2x_4+x_0^2x_1^3x_4+x_0^2x_1^2x_3x_4+x_0^2x_1^2x_4^2+x_0^2x_1x_2^2x_3+x_0^2x_1x_2x_3^2+x_0^2x_1x_3^3+x_0^2x_1x_3x_4^2+x_0^2x_1x_4^3+x_0^2x_2^2x_4^2+x_0^2x_3^2x_4^2+x_0x_1^3x_4^2+x_0x_1^2x_3^2x_4+x_0x_1^2x_4^3+x_0x_1x_2^2x_3^2+x_0x_1x_3^4+x_1^3x_3^2x_4+x_1^3x_3x_4^2$.
		
		
	\end{proposition}

	\begin{proof}
	
Put $T=\mathbb F_{2}[x_{0},x_{1},x_{2},x_{3},x_{4}], \qquad S=T^{G}, \qquad A=\mathbb F_{2}[f_{1},\ldots,f_{9}]\subseteq S$. The action of the generators of \(G=D_{8}\) on the degree-one component \(T_{1}=\Span_{\mathbb F_2}\{x_0,x_1,x_2,x_3,x_4\}\) is
			\[
			\begin{aligned}
				&r(x_{0})=x_{0},\qquad
				r(x_{1})=x_{0}+x_{1},\qquad
				r(x_{2})=x_{1}+x_{2},\\
				&r(x_{3})=x_{3},\qquad
				r(x_{4})=x_{3}+x_{4},
			\end{aligned}
			\]
			and
			\[
			\begin{aligned}
				&s(x_{0})=x_{0},\qquad
				s(x_{1})=x_{0}+x_{1},\qquad
				s(x_{2})=x_{0}+x_{2},\\
				&s(x_{3})=x_{3},\qquad
				s(x_{4})=x_{0}+x_{3}+x_{4}.
			\end{aligned}
			\]
Direct substitution gives $r(f_i)=f_i=s(f_i) \qquad (1\leq i\leq 9)$, so \(A\subseteq S\). We first exhibit a homogeneous system of parameters for \(S\). Set
			\[
			p_{1}=f_{1}=x_{0},\qquad
			p_{2}=f_{2}=x_{3},\qquad
			p_{3}=f_{3}=x_{0}x_{1}+x_{1}^{2},
			\qquad
			p_{4}=f_{6},\qquad
			p_{5}=f_{7}.
			\]
			Their degrees are $1,1,2,4,4$. We claim that their only common zero over \(\overline{\mathbb F}_{2}\) is the origin. Indeed, if $p_{1}=p_{2}=p_{3}=p_{4}=p_{5}=0$, then \(p_{1}=p_{2}=0\) gives $x_{0}=x_{3}=0$. It follows from $p_{3}=x_{0}x_{1}+x_{1}^{2}$ that $x_{1}^{2}=0$, and hence \(x_{1}=0\). Under the substitutions
			\(x_{0}=x_{1}=x_{3}=0\), the two remaining parameters reduce to $p_{4}=f_{6}=x_{4}^{4}, \qquad p_{5}=f_{7}=x_{2}^{4}$. Thus \(x_{4}=x_{2}=0\), proving the claim. Consequently,
			\[
			T/(p_{1},p_{2},p_{3},p_{4},p_{5})T
			\]
			is finite-dimensional over \(\mathbb F_{2}\). Hence \(T\) is module-finite over $P=\mathbb F_{2}[p_{1},p_{2},p_{3},p_{4},p_{5}]$. Since \(P\subseteq S\subseteq T\) and \(P\) is Noetherian, \(S\) is also
			module-finite over \(P\). Therefore \(p_{1},\ldots,p_{5}\) form a homogeneous system of parameters for \(S\). By Symonds' bound for secondary invariants, \(S\) is generated as a
			\(P\)-module by homogeneous elements of degree at most
			\[
			(1-1)+(1-1)+(2-1)+(4-1)+(4-1)=7.
			\]
			It remains to verify that
			\[
			A_d=S_d\qquad(0\leq d\leq7).
			\]
			For each \(d\geq0\), let
			\[
			\mathcal M_d=
			\left\{x_0^{a_0}\cdots x_4^{a_4}:a_0+\cdots+a_4=d\right\}
			\]
			be the standard monomial basis of \(T_d\), and put
			\[
			N_d=|\mathcal M_d|=\binom{d+4}{4}.
			\]
			Let \(R_d\) and \(S_d'\) denote the matrices, with respect to
			\(\mathcal M_d\), of the transformations induced by \(r\) and \(s\),
			respectively, and define
			\[
			\Phi_d=
			\begin{pmatrix}
			R_d-I\\ S_d'-I
			\end{pmatrix}.
			\]
			Then
			\[
			S_d=\ker(R_d-I)\cap\ker(S_d'-I)=\ker\Phi_d,
			\]
			and hence
			\[
			\dim_{\mathbb F_2}S_d=N_d-\operatorname{rank}_{\mathbb F_2}\Phi_d.
			\]
			Let \(C_d\) be the matrix whose columns are the coefficient vectors,
			in the basis \(\mathcal M_d\), of all products
			\(f_1^{a_1}\cdots f_9^{a_9}\) satisfying
			\[
			a_{1}+a_{2}+2a_{3}+3a_{4}
			+4(a_{5}+a_{6}+a_{7}+a_{8})+6a_{9}=d.
			\]
			The column space of \(C_{d}\) is precisely \(A_{d}\). Exact Gaussian elimination over \(\mathbb F_{2}\) gives
			\[
			\begin{array}{c|rrrrrrrr}
				d&0&1&2&3&4&5&6&7\\ \hline
				N_{d}
				&1&5&15&35&70&126&210&330\\
				\operatorname{rank}\Phi_{d}
				&0&3&11&28&55&103&173&277\\
				\dim_{\mathbb F_{2}}S_{d}
				&1&2&4&7&15&23&37&53\\
				\operatorname{rank}C_{d}
				&1&2&4&7&15&23&37&53.
			\end{array}
			\]
			Therefore $\dim_{\mathbb F_{2}}A_{d} =\dim_{\mathbb F_{2}}S_{d} \qquad(0\leq d\leq7)$. Since \(A_{d}\subseteq S_{d}\), it follows that $A_{d}=S_{d} \qquad(0\leq d\leq7)$. Every homogeneous \(P\)-module generator of \(S\) has degree at most \(7\), and hence belongs to \(A\). Since \(P\subseteq A\), we conclude that $ S=P\cdot S_{\leq7}\subseteq A$. The reverse inclusion \(A\subseteq S\) has already been proved, so $S=A=\mathbb F_{2}[f_{1},\ldots,f_{9}]$. Finally, we verify minimality. Let $S_{+}=\bigoplus_{d>0}S_{d}$. Since \(S=A\), the space \((S_{+}^{2})_{d}\) is spanned by the products $f_{1}^{a_{1}}\cdots f_{9}^{a_{9}}$ of weighted degree \(d\) for which $a_{1}+\cdots+a_{9}\geq2$. A further exact row reduction gives
			\[
			\begin{array}{c|rrrrrrr}
				d&1&2&3&4&5&6&7\\ \hline
				\dim_{\mathbb F_{2}}S_{d}
				&2&4&7&15&23&37&53\\
				\dim_{\mathbb F_{2}}(S_{+}^{2})_{d}
				&0&3&6&11&23&36&53\\ \hline
				\dim_{\mathbb F_{2}}
				\bigl(S_{d}/(S_{+}^{2})_{d}\bigr)
				&2&1&1&4&0&1&0.
			\end{array}
			\]
			The corresponding nonzero indecomposable quotients have bases
			\[
			\begin{array}{c|c}
				d&
				\text{basis of }S_{d}/(S_{+}^{2})_{d}\\ \hline
				1&\overline f_{1},\overline f_{2}\\
				2&\overline f_{3}\\
				3&\overline f_{4}\\
				4&\overline f_{5},\overline f_{6},
				\overline f_{7},\overline f_{8}\\
				6&\overline f_{9}.
			\end{array}
			\]
			Thus none of \(f_{1},\ldots,f_{9}\) can be omitted. Hence they form a minimal homogeneous generating set of \(S\), with degrees $1,1,2,3,4,4,4,4,6$. In particular, $\beta\!\left(\mathbb F_{2}[U]^{G}\right)=6$.
		
	\end{proof}

\end{document}
